%%%%%%%%%%%%%%%%%%%%%%%%%%%%%%%%%%%%%%%%%%%%%%%%%%%%%%%%%%%%%%%%%%%%%%%%%%%%%
%% 03_arch_clock_concept.tex
%% Section: Clock Concept
%% Chapter: Architecture Concepts
%%
%% Describes the Clock Generation Unit (CGU) shared by both chip variants
%% and the concrete clock-frequency requirements for Variant~A
%% (RTL top RV_NoPipeline).
%%
%% Integration in 00_rwurv64i_all.tex:
%%   \section{Clock Concept}
%%   \input{\TEXTarch 03_arch_history_tab04}
%%   \input{\TEXTarch 03_arch_clock_concept}
%%
%% Requires packages: tabularx, booktabs, float
%%
%% v0.1  2026-07-12  Claude-AI  Initial version -- RV_NoPipeline clock requirements
%%%%%%%%%%%%%%%%%%%%%%%%%%%%%%%%%%%%%%%%%%%%%%%%%%%%%%%%%%%%%%%%%%%%%%%%%%%%%

%%%%%%%%%%%%%%%%%%%%%%%%%%%%%%%%%%%%%%%%%%%%%%%%%%%%%%%%%%%%%%%%%%%%%%%%%%%%%
\subsection{Clock Generation Unit (CGU) Overview}
\label{subsec:cguOverview}
%%%%%%%%%%%%%%%%%%%%%%%%%%%%%%%%%%%%%%%%%%%%%%%%%%%%%%%%%%%%%%%%%%%%%%%%%%%%%

Both chip variants (Table~\ref{tab:asChipVariants}) share the same Clock
Generation Unit IP (\texttt{as\_cgu\_top} / \texttt{as\_cgucore}).
The chip has a single external clock input, \texttt{clk\_i}, from which the
CGU derives all internal clock domains using free-running counter-based
dividers.
There is no PLL: every derived clock is an integer division of \texttt{clk\_i}.

\begin{table}[H]
\caption{CGU Clock Domains}
\label{tab:asCguDomains}
\centering
\begin{tabularx}{\textwidth}{|l |l |X |X|}
  \hline
  Output & Divider (fixed) & Consumers & Notes \\
  \hline
  \hline
  \texttt{clk\_core\_o} & \texttt{clk\_core\_div} = 80 &
    CPU core, GPIO, UART0, D-Mem/Scratchpad, D-/I-Cache (Variant~B) &
    Functional clock for the CPU and almost all peripherals \\
  \hline
  \texttt{clk\_bus1\_o} & \texttt{clk\_bus1\_div} = 80 & Reserved & Not currently consumed in \texttt{as\_top\_mem} \\
  \hline
  \texttt{clk\_bus2\_o} & \texttt{clk\_bus2\_div} = 100 & Reserved & Not currently consumed in \texttt{as\_top\_mem} \\
  \hline
  \texttt{clk\_qspi\_o} & \texttt{clk\_qspi\_div} = 4 & QSPI controller (Variant~B only) &
    Not present in Variant~A (no QSPI controller) \\
  \hline
  \texttt{tck\_i} & --- (external JTAG clock) & JTAG TAP, I-Mem scan chain, boundary scan &
    Fully asynchronous to \texttt{clk\_i}; no timing relationship is analysed
    between the two domains \\
  \hline
\end{tabularx}
\end{table}

\textbf{The divider ratios are fixed at RTL compile time, not run-time
configurable.}
Unlike the QSPI and UART peripherals (which expose a genuine \texttt{CLKDIV}
Wishbone register), the CGU implements only a read/write peripheral
\textbf{ID register}; the four divider ratios above are SystemVerilog
\texttt{localparam}s in \texttt{as\_pack.sv} and take effect unconditionally.
Changing a divider therefore requires an RTL/constant change and a full
re-synthesis -- it cannot be adjusted by software or by writing a register.

\subsubsection{Scan-Test Clock Bypass}

During normal operation the CPU and Wishbone peripherals are clocked by
\texttt{clk\_div\_s}, which is generated in \texttt{as\_top\_mem} as:
\begin{verbatim}
assign clk_div_s = dr_cap_s ? clk_i : clk_core_s;
\end{verbatim}
\texttt{dr\_cap\_s} is driven by the JTAG TAP.
In normal function mode (\texttt{dr\_cap\_s}\,=\,0) the CPU runs on the
divided \texttt{clk\_core\_s}.
In scan-test mode (\texttt{dr\_cap\_s}\,=\,1) the divider is bypassed and the
CPU is clocked directly by the raw external \texttt{clk\_i}, so that an ATE
tester can single-step or free-run the scan chain independently of the CGU
divide ratio.
The CGU itself (\texttt{as\_cgu\_top}) always runs on the undivided
\texttt{clk\_i}; only its divided outputs are subject to the bypass mux.

%%%%%%%%%%%%%%%%%%%%%%%%%%%%%%%%%%%%%%%%%%%%%%%%%%%%%%%%%%%%%%%%%%%%%%%%%%%%%
\subsection{Clock Requirements -- Variant A (RV\_NoPipeline)}
\label{subsec:clockReqNoPipeline}
%%%%%%%%%%%%%%%%%%%%%%%%%%%%%%%%%%%%%%%%%%%%%%%%%%%%%%%%%%%%%%%%%%%%%%%%%%%%%

\subsubsection{Status}

No X-FAB XH035 gate-level static timing analysis (STA) has been run yet for
\texttt{RV\_NoPipeline}.
The Genus synthesis script (\texttt{scripts/01b\_rv64\_syn.tcl}) still carries
a placeholder constraint (\texttt{create\_clock -period 20.000}, i.e.\ 50\,MHz
on \texttt{clk\_i}) and is explicitly marked in its own comments as
provisional, pending SRAM-macro integration for I-Mem/D-Mem.
The recommendation below is therefore a conservative pre-STA target, not a
timing-signed-off number.

\subsubsection{Recommended Frequencies}

\begin{table}[H]
\caption{RV\_NoPipeline -- Recommended Clock Frequencies}
\label{tab:asNoPipelineClk}
\centering
\begin{tabularx}{\textwidth}{|X |X |X|}
  \hline
  Signal & Recommended value & Rationale \\
  \hline
  \hline
  \texttt{clk\_i} (external) & 100\,MHz (10\,ns period) &
    Matches the operating point already exercised by every existing xsim
    and cocotb testbench (\texttt{clk\_2\_t}\,=\,5\,ns half-period) \\
  \hline
  \texttt{clk\_core\_s} (CPU core, \texttt{clk\_i}/80) & 1.25\,MHz (800\,ns period) &
    Matches \texttt{clk\_80\_t}\,=\,400\,ns in the testbenches and the
    1.25\,MHz core clock already assumed in the UART timing calculation
    (Section~\ref{subsec:uartsimtest}) \\
  \hline
  \texttt{tck\_i} (JTAG) & $\leq$ \texttt{clk\_i}/4 (typ.\ 25\,MHz or less) &
    Conservative default; JTAG is asynchronous to \texttt{clk\_i}
    (Table~\ref{tab:asCguDomains}) and not on the CPU's timing-critical path \\
  \hline
\end{tabularx}
\end{table}

\subsubsection{Rationale}

\begin{itemize}
  \item \textbf{Why 1.25\,MHz and not higher:} \texttt{RV\_NoPipeline} is a
    multi-cycle, non-pipelined core (Section~\ref{sec:archcpu}): each
    instruction's critical path runs through instruction decode, register
    file, the ALU, and (for loads/stores) the Wishbone master/slave BPI
    handshake, all within states of a single clock.
    X-FAB XH035 is a 350\,nm process; without a completed STA on the actual
    gate-level netlist (SRAM macros not yet integrated, see Status above),
    committing to a high frequency is premature.
    1.25\,MHz leaves a very large safety margin against any plausible
    critical-path length on a 350\,nm standard-cell library.
  \item \textbf{Why keep the already-validated operating point:} every xsim
    and cocotb regression test (58 xsim / 57 cocotb tests, see
    Section~\ref{sec:simtest01}) already runs at \texttt{clk\_i}\,=\,100\,MHz
    / \texttt{clk\_core\_s}\,=\,1.25\,MHz, and the UART bit-timing
    calculation in Section~\ref{subsec:uartsimtest} is derived from exactly
    this core clock.
    Choosing the same frequency for the first tapeout avoids re-deriving or
    re-validating any functional timing (UART baud rate, GPIO polling loops,
    etc.).
  \item \textbf{Fixed divider consequence:} because \texttt{clk\_core\_div}\,=\,80
    is fixed at compile time (Section~\ref{subsec:cguOverview}), the external
    \texttt{clk\_i} frequency is always exactly $80\times$ the desired core
    frequency.
    There is no register to trade off external frequency against core
    frequency at run time.
\end{itemize}

\subsubsection{Open Items}

\textbf{Action (CLK-01):} Once I-Mem/D-Mem SRAM macros are integrated into
the \texttt{RV\_NoPipeline} netlist, run Genus \texttt{report\_timing} with a
proper \texttt{create\_generated\_clock -divide\_by 80} constraint on
\texttt{clk\_core\_o} (in addition to the existing constraint on
\texttt{clk\_i}, which today only reflects the scan-test bypass path) and
confirm that 1.25\,MHz core / 100\,MHz external clears timing with adequate
margin at the worst-case corner.
If margin is large, a higher core frequency may be considered for a later
revision -- this requires changing the \texttt{clk\_core\_div} localparam in
\texttt{as\_pack.sv} and re-synthesising, not a register write.
